\scanpage{183}
\mspage{I bis}
\margin{Mai 1986\\Juillet 1986}
$\unsure{\mathcal{A}|A}$ alg. d'Az. de rang $4$ sur $k$, on veut expliciter sa structure multiplicative en termes de la forme $\det$, de l'élément $1$, et de la structure \uline{spéciale} orthogonale $e\in\bigwedge^4\mathcal{A}$ can. \deleted{\unsure{associée}} associée à $\det$ ($e\otimes e=\operatorname{dis}(\det)$).
Comme on connaît la forme $\operatorname{Tr}xy$ par la formule
\margin{\uline{NB} $\operatorname{Tr}$ connu en termes de $\det$ et de $1$, comme $\operatorname{Tr}(x)=\Psi_{\det}(x,1)$ i.e.}
\[
\det(x+y)=\det x+\det y+\operatorname{Tr}x\operatorname{Tr}y-\operatorname{Tr}xy
\]
\[
\Psi_{\det\deleted{\unsure{(xy)}}}(x,y)=\operatorname{Tr}x\operatorname{Tr}y-\underbrace{\operatorname{Tr}xy}_{\varphi(x,y)},
\]
\deleted{pour} et que celle-ci est non dég., pour connaître $xy$, il suffit de connaître
\[
\operatorname{Tr}xyz\qquad\forall z\in\mathcal{A}.
\]
Donc je cherche une formule pour $\operatorname{Tr}xyz$. Comme $\operatorname{Tr}(xy)z$ et $\operatorname{Tr}(yx)z$ ne sont pas égaux en général, i.e la structure mult\inserted{de $\mathcal{A}$} n'est pas commutative \deleted{(donc} On voit que $\det,1$ sans $e$ ne suffisent à expliciter $\operatorname{Tr}xyz$ (car ils ne déterminent la structure multiplicative que mod passage à l'inverse, i.e. mod l'involution
\[
x\longmapsto x'=\operatorname{Tr}(x).1-x.
\]
On note aussi que $\operatorname{Tr}xyz$ ne dépend que du triple $\{x,y,z\}$ et d'un ordre circulaire sur l'\unsure{ens.} des trois termes. On a, \deleted{pour} si
\[
\begin{aligned}
\tau&=\operatorname{Tr}xyz\quad(=\operatorname{Tr}yzx=\operatorname{Tr}zxy)\\
\tau'&=\operatorname{Tr}(zyx)\bigl(=\operatorname{Tr}(yxz)=\operatorname{Tr}(xzy)\bigr)
\end{aligned}
\]
\margin{182}
